\documentclass[12pt, titlepage]{article}

\usepackage{amsmath, mathtools}

\usepackage[round]{natbib}
\usepackage{amsfonts}
\usepackage{amssymb}
\usepackage{graphicx}
\usepackage{colortbl}
\usepackage{xr}
\usepackage{hyperref}
\usepackage{longtable}
\usepackage{xfrac}
\usepackage{tabularx}
\usepackage{float}
\usepackage{siunitx}
\usepackage{booktabs}
\usepackage{multirow}
\usepackage[section]{placeins}
\usepackage{caption}
\usepackage{fullpage}
\usepackage{comment}

\hypersetup{
bookmarks=true,     % show bookmarks bar?
colorlinks=true,       % false: boxed links; true: colored links
linkcolor=red,          % color of internal links (change box color with linkbordercolor)
citecolor=blue,      % color of links to bibliography
filecolor=magenta,  % color of file links
urlcolor=cyan          % color of external links
}

\usepackage{array}

\input{../../Comments}
%% Common Parts

\newcommand{\progname}{Software Engineering} % PUT YOUR PROGRAM NAME HERE
\newcommand{\authname}{Team 8, AlgoCatan
\\ Jake Read
\\ Rebecca Di Filippo
\\ Matthew Cheung
\\ Sunny Yao} % AUTHOR NAMES

\usepackage{hyperref}
    \hypersetup{colorlinks=true, linkcolor=blue, citecolor=blue, filecolor=blue,
                urlcolor=blue, unicode=false}
    \urlstyle{same}
                                


\begin{document}

\title{Module Interface Specification for \progname{}}

\author{\authname}

\date{\today}

\maketitle

\pagenumbering{roman}

\section{Revision History}

\begin{tabularx}{\textwidth}{p{3cm}p{2cm}X}
\toprule {\bf Date} & {\bf Version} & {\bf Notes}\\
\midrule
11/13/2025 & 1.0 & Draft Rev Minus 1\\
01/17/2026 & 1.1 & Draft Rev 0\\
04/06/2026 & 1.2 & Revision for Final Doc\\
\bottomrule
\end{tabularx}

~\newpage

\section{Symbols, Abbreviations and Acronyms}

See SRS Documentation at \url{https://github.com/AlgoCatan/RLCatan}.


%\wss{Also add any additional symbols, abbreviations or acronyms}

\newpage

\tableofcontents

\newpage

\pagenumbering{arabic}

\section{Introduction}

The following document details the Module Interface Specifications for
our project, AlgoCatan. The final implemented system is a web-based arena for
playing against trained \emph{Catan} bots, replaying past games, viewing
strategic analysis, and requesting natural-language move explanations.
Its architecture combines a training pipeline, PPO-based AI models, an Elo
league, a curriculum learning manager, a backend API server, persistent game
storage, and a React frontend.

Complementary documents include the Software Requirements Specification
and Module Guide. The full documentation and implementation can be
found at \url{https://github.com/AlgoCatan/RLCatan}.

\section{Notation}

%\wss{You should describe your notation.  You can use what is below as
%  a starting point.}

The structure of the MIS for modules comes from \citet{HoffmanAndStrooper1995},
with the addition that template modules have been adapted from
\cite{GhezziEtAl2003}.  The mathematical notation comes from Chapter 3 of
\citet{HoffmanAndStrooper1995}.  For instance, the symbol := is used for a
multiple assignment statement and conditional rules follow the form $(c_1
\Rightarrow r_1 | c_2 \Rightarrow r_2 | ... | c_n \Rightarrow r_n )$.

The following table summarizes the primitive data types used by \progname. 

\begin{center}
  \renewcommand{\arraystretch}{1.2}
  \noindent 
  \begin{tabular}{l l p{7.5cm}} 
  \toprule 
  \textbf{Type/Symbol} & \textbf{Notation} & \textbf{Description}\\ 
  \midrule
  % Primitive Types
  boolean & $\mathbb{B}$ & a logical value, either true ($\top$) or false ($\bot$) \\
  character & char & a single symbol or digit\\
  integer & $\mathbb{Z}$ & a number without a fractional component in (-$\infty$, $\infty$) \\
  natural number & $\mathbb{N}$ & a number without a fractional component in [1, $\infty$) \\
  real & $\mathbb{R}$ & any number in (-$\infty$, $\infty$)\\
  string & $\Sigma^*$ & a sequence of zero or more characters. \\
  sequence & seq of T & a list of zero or more elements of the same type T. \\
  tuple & (T1, T2, ...) & a finite, ordered list of elements, potentially of
                          different types. \\
  dictionary & dict & a collection of key-value mappings (ex, KeyType $\rightarrow$ ValueType). \\
  function & T1 $\rightarrow$ T2 & a function mapping an input of type T1 to
                          an output of type T2. \\
  true & $\top$ & the boolean value for true. \\
  false / null & $\bot$ & the boolean value for false, or a value that is
                          undefined, null, or void (no output). \\
  \bottomrule
  \end{tabular} 
  \end{center}

\noindent
In addition, \progname \ uses abstract
 data types (e.g., \texttt{FrameData}, 
 \texttt{GameStateData}, \texttt{AIMove})
  which are defined by their use in the 
  modules that hide their implementation 
  details. Local functions are described by
   giving their type signature followed by
    their specification.

\section{Module Decomposition}

The following table is taken directly from the Module Guide document for this project.

\begin{table}[h!]
\centering
\begin{tabular}{p{0.3\textwidth} p{0.6\textwidth}}
\toprule
\textbf{Level 1} & \textbf{Level 2}\\
\midrule

\multirow{2}{0.3\textwidth}{Hardware-Hiding Module}
& Hardware-Hiding Module (OS) (Deferred) \\
& Computer Vision Model (Deferred) \\
\midrule

\multirow{3}{0.3\textwidth}{Behaviour-Hiding Module}
& User Interface \\
& Game State Manager \\
& Training Pipeline \\
\midrule

\multirow{6}{0.3\textwidth}{Software Decision Module}
& AI Model \\
& Game State Database \\
& Backend API Server \\
& Elo League System \\
& Curriculum Learning Manager \\
& Explanation Pipeline \\
\bottomrule

\end{tabular}
\end{table}

\newpage
~\newpage


%%%%%%%%%%%%%%%%%ADDED MODULES BELOW THIS LINE%%%%%%%%%%%%%%%%%%%%%
\section{MIS of Hardware-Hiding Module (Deferred)} \label{M1}

\subsection{Module}
Hardware-Hiding Module (M1).
This module is retained for future scope only and is not part of the current implemented runtime path.

\subsection{Uses}
\begin{itemize}
    \item \textbf{FrameData (Data Type):} Represents the raw image data captured from the camera or a simulated source.
    \item \textbf{HardwareInitError, CaptureError, HardwareShutdownError (Exception Types):} Indicate errors that may occur during hardware initialization, frame capture, or shutdown processes.
\end{itemize}

\subsection{Syntax}

\subsubsection{Exported Constants}
\begin{itemize}
\item \texttt{CAMERA\_RESOLUTION}: tuple of $(\mathbb{Z}, \mathbb{Z})$ - width and height of captured frames
\item \texttt{FRAME\_RATE}: $\mathbb{R}$ - frames per second captured by the hardware layer
\end{itemize}

\subsubsection{Exported Access Programs}
\begin{center}
\begin{tabular}{p{5cm} p{3cm} p{3cm} p{5cm}}
\hline
\textbf{Name} & \textbf{In} & \textbf{Out} & \textbf{Exceptions} \\
\hline
initializeHardware & - & $\mathbb{B}$ & HardwareInitError \\
captureFrame & - & FrameData & CaptureError \\
shutdownHardware & - & $\mathbb{B}$ & HardwareShutdownError \\
isInitialized & - & $\mathbb{B}$ & - \\
\hline
\end{tabular}
\end{center}

\subsection{Semantics}

\subsubsection{State Variables}
\begin{itemize}
\item \texttt{hardwareStatus}: $\mathbb{B}$ \\
      Indicates whether the hardware subsystem has been successfully initialized.
\item \texttt{frameBuffer}: FrameData \\
      Stores the most recent frame captured from the camera.
\end{itemize}

\subsubsection{Environment Variables}
\begin{itemize}
\item \texttt{CameraDevice}: external hardware input for image capture
\item \texttt{DisplayInterface}: external interface for visualization
\end{itemize}

\subsubsection{Assumptions}
\begin{itemize}
  \item The CameraDevice is available and functional within the operating environment.
  \item Calling modules will invoke \texttt{isInitialized()} before executing 
        \texttt{captureFrame()} or \texttt{shutdownHardware()}.
\end{itemize}

\subsubsection{Access Routine Semantics}

\noindent initializeHardware():
\begin{itemize}
\item transition: Attempts connection to \texttt{CameraDevice}.  
      If successful: \texttt{hardwareStatus} := $\top$.
\item output: Returns initialization success status: $\mathbb{B}$
\item exception: raises \texttt{HardwareInitError} if initialization fails.
\end{itemize}

\noindent captureFrame():
\begin{itemize}
\item transition: Reads a new frame from \texttt{CameraDevice} into \texttt{frameBuffer}.
\item output: Returns the current \texttt{frameBuffer} of type FrameData.
\item exception: raises \texttt{CaptureError} if \texttt{hardwareStatus} = $\bot$
                 or frame capture fails.
\end{itemize}

\noindent shutdownHardware():
\begin{itemize}
\item transition: Releases device resources and sets  
      \texttt{hardwareStatus} := $\bot$.
\item output: Returns shutdown success status: $\mathbb{B}$
\item exception: raises \texttt{HardwareShutdownError} if hardware cleanup fails.
\end{itemize}

\noindent isInitialized():
\begin{itemize}
\item transition: -
\item output: returns \texttt{hardwareStatus}: $\mathbb{B}$
\item exception: -
\end{itemize}

\subsubsection{Local Functions}
\begin{itemize}
\item \texttt{checkDeviceConnection()}: verifies device presence
\item \texttt{allocateBufferMemory()}: allocates buffer for captured frames
\item \texttt{releaseResources()}: frees hardware-related resources
\end{itemize}


%%%%%%%%%%%%%%%%%%%%%%%%%%%%%%%%%%%%%%%%%%%%%%

\section{MIS of Computer Vision Module (Deferred)} \label{M2}

\subsection{Module}
Computer Vision Module (M2).
This module is retained for future scope only and is not part of the current implemented runtime path.

\subsection{Uses}
\begin{itemize}
    \item \textbf{M1.captureFrame():} Provides a single frame of image data from the hardware or simulated camera.
    \item \textbf{FrameData (Data Type):} Represents the raw image frame used for analysis.
    \item \textbf{GameStateData (Data Type):} Abstract representation of detected game state.
    \item \textbf{ProcessingError, DetectionError (Exception Types):} Raised during processing or detection failures.
    \item \textbf{seq of Elements, dict} (Derived Data Types):  
          Sequences and mappings used to store detected features and confidence metrics.
\end{itemize}

\subsection{Syntax}

\subsubsection{Exported Constants}
\begin{itemize}
\item \texttt{DETECTION\_THRESHOLD}: $\mathbb{R}$ - minimum required confidence for object recognition
\item \texttt{MODEL\_PATH}: $\Sigma^{*}$ - file system path to trained CV model
\end{itemize}

\subsubsection{Exported Access Programs}
\begin{center}
\begin{tabular}{p{5cm} p{4cm} p{4cm} p{3cm}}
\hline
\textbf{Name} & \textbf{In} & \textbf{Out} & \textbf{Exceptions} \\
\hline
processFrame & frame:FrameData & GameStateData & ProcessingError \\
detectBoardElements & frame:FrameData & seq of Elements & DetectionError \\
getConfidenceMetrics & - & dict & - \\
\hline
\end{tabular}
\end{center}

\subsection{Semantics}

\subsubsection{State Variables}
\begin{itemize}
\item \texttt{lastFrame}: FrameData \\
      Stores the most recent processed frame.
\item \texttt{lastConfidence}: dict \\
      A mapping from element identifiers ($\Sigma^{*}$) to confidence scores ($\mathbb{R}$).
\end{itemize}

\subsubsection{Environment Variables}
None.

\subsubsection{Assumptions}
\begin{itemize}
    \item All \texttt{FrameData} inputs are valid and camera-calibrated.
    \item The model at \texttt{MODEL\_PATH} is present and loaded into memory before use.
\end{itemize}

\subsubsection{Access Routine Semantics}

\noindent processFrame(frame):
\begin{itemize}
\item transition: Executes the full CV pipeline on $frame$;  
      updates \texttt{lastFrame} := frame.
\item output: Returns detected game state as \texttt{GameStateData}.
\item exception: raises \texttt{ProcessingError} if any stage of the pipeline fails.
\end{itemize}

\noindent detectBoardElements(frame):
\begin{itemize}
\item transition: Runs object detection model on $frame$;  
      updates \texttt{lastFrame} := frame and updates \texttt{lastConfidence}.
\item output: Returns $seq$ of detected Elements.
\item exception: raises \texttt{DetectionError} if detection fails.
\end{itemize}

\noindent getConfidenceMetrics():
\begin{itemize}
\item transition: -
\item output: Returns the \texttt{lastConfidence} mapping (a dict from $\Sigma^{*}$ to $\mathbb{R}$).
\item exception: -
\end{itemize}

\subsubsection{Local Functions}
\begin{itemize}
\item \texttt{calibrateCamera()}: Corrects lens distortion in input frames.
\item \texttt{filterNoise()}: Removes low-confidence or spurious detections.
\item \texttt{parseElementsToState()}: Converts detected Elements into \texttt{GameStateData}.
\end{itemize}


%%%%%%%%%%%%%%%%%%%%%%%%%%%%%%%%%%%%%%%%%%%%%%%

\section{MIS of User Interface Module} \label{M3}

\subsection{Module}
User Interface (M3)

\subsection{Uses}

\begin{itemize}
    \item \textbf{GameStateData}: Abstract representation of the current game state.
    \item \textbf{AIMove}: Encoded AI-generated move.
    \item \textbf{seq of Corrections}: A sequence of user-provided corrections to the game state.
    \item \textbf{RenderError}: Exception raised when visualization or rendering fails.
\end{itemize}

\subsection{Syntax}

\subsubsection{Exported Constants}
\begin{itemize}
\item \texttt{REFRESH\_RATE}: $\mathbb{R}$ - frequency of UI updates (Hz)
\end{itemize}

\subsubsection{Exported Access Programs}
\begin{center}
\begin{tabular}{p{5cm} p{4cm} p{4cm} p{3cm}}
\hline
\textbf{Name} & \textbf{In} & \textbf{Out} & \textbf{Exceptions} \\
\hline
renderBoard & state:GameStateData & - & RenderError \\
displayAIMove & move:AIMove & - & - \\
getUserCorrections & - & seq of Corrections & - \\
\hline
\end{tabular}
\end{center}

\subsection{Semantics}

\subsubsection{State Variables}
\begin{itemize}
\item \texttt{uiState}: Abstract UI representation (type unspecified, internal)
\item \texttt{correctionQueue}: seq of Corrections
      \quad A sequence of user-provided corrections where each correction may include:
      \begin{itemize}
          \item element identifier: $\Sigma^{*}$
          \item correction value: arbitrary derived type
      \end{itemize}
\end{itemize}

\subsubsection{Environment Variables}
\begin{itemize}
\item \texttt{window}: graphical display device managed by the environment
\end{itemize}

\subsubsection{Assumptions}
\begin{itemize}
    \item The \texttt{window} device is functional and compatible with the UI rendering framework.
    \item User interactions are delivered to the system via environment-level event handling.
\end{itemize}

\subsubsection{Access Routine Semantics}

\noindent renderBoard(state):
\begin{itemize}
\item transition: Updates \texttt{uiState} := graphical representation of $state$;
      renders \texttt{uiState} onto \texttt{window}.
\item output: -
\item exception: raises \texttt{RenderError} if rendering fails.
\end{itemize}

\noindent displayAIMove(move):
\begin{itemize}
\item transition: Updates \texttt{uiState} to include the visualization of $move$;
      refreshes \texttt{window}.
\item output: -
\item exception: -
\end{itemize}

\noindent getUserCorrections():
\begin{itemize}
\item transition: \texttt{tmp} := \texttt{correctionQueue};
      \texttt{correctionQueue} := empty sequence.
\item output: Returns \texttt{tmp}: seq of Corrections.
\item exception: -
\end{itemize}

\subsubsection{Local Functions}
\begin{itemize}
\item \texttt{updateDOM()}: Updates internal UI representation elements.
\item \texttt{highlightElements()}: Highlights tiles, pieces, or game move indicators.
\item \texttt{onUserInput()}: Adds user-generated corrections to \texttt{correctionQueue}.
\end{itemize}

%%%%%%%%%%%%%%%%%%%%%%%%%%%%%%%%%%%%%%%%%%%%%%

\section{MIS of Game State Manager Module} \label{M4}

\subsection{Module}
Game State Manager (M4)

\subsection{Uses}

\begin{itemize}
    \item \textbf{MoveData}: Represents the details of a player's move, including action type and parameters.
    \item \textbf{GameStateData}: Abstract representation of the complete game state.
    \item \textbf{InvalidMoveError}: Raised when an illegal move is applied.
    \item \textbf{M7.writeState()}: Persists updated game state to storage.
\end{itemize}

\subsection{Syntax}

\subsubsection{Exported Constants}
\begin{itemize}
\item \texttt{MAX\_PLAYERS}: $\mathbb{Z}^{+}$ = 4
\end{itemize}

\subsubsection{Exported Access Programs}
\begin{center}
\begin{tabular}{p{5cm} p{4cm} p{4cm} p{3cm}}
\hline
\textbf{Name} & \textbf{In} & \textbf{Out} & \textbf{Exceptions} \\
\hline
updateState & move:MoveData & - & InvalidMoveError \\
getState & - & GameStateData & - \\
validateMove & move:MoveData & $\mathbb{B}$ & - \\
\hline
\end{tabular}
\end{center}

\subsection{Semantics}

\subsubsection{State Variables}
\begin{itemize}
\item \texttt{currentGameState}: GameStateData \\
      The full internal representation of the game board.

\item \texttt{playerAssets}: dict \\
      Mapping from player identifiers ($\Sigma^{*}$) to their resource, structure,  
      and score information. Each value may include:
      \begin{itemize}
          \item resources: dict from resource type ($\Sigma^{*}$) to quantity ($\mathbb{Z}$)
          \item structures: seq of structure types
          \item score: $\mathbb{Z}_{\ge 0}$
      \end{itemize}
\end{itemize}

\subsubsection{Environment Variables}
None.

\subsubsection{Assumptions}
\begin{itemize}
    \item All instances of \texttt{MoveData} conform to the expected internal format.
    \item Calling modules will evaluate \texttt{validateMove(move)} before invoking \texttt{updateState(move)}.
\end{itemize}

\subsubsection{Access Routine Semantics}

\noindent updateState(move):
\begin{itemize}
  \item transition:
  \[
  (
    \texttt{validateMove(move)} = \top \Rightarrow
      \text{\texttt{currentGameState}} := \text{\texttt{applyMove}} (\text{move})
  |
    \texttt{validateMove(move)} = \bot \Rightarrow
      \text{-}
  )
  \]
\item output: -
\item exception: raises \texttt{InvalidMoveError} if \texttt{validateMove(move)} = $\bot$.
\end{itemize}

\noindent getState():
\begin{itemize}
\item transition: -
\item output: Returns a copy of \texttt{currentGameState}: GameStateData.
\item exception: -
\end{itemize}

\noindent validateMove(move):
\begin{itemize}
\item transition: -
\item output: $\mathbb{B}$ - $\top$ iff $move$ is valid according to Catan rules and the
      current values of \texttt{currentGameState} and \texttt{playerAssets}.
\item exception: -
\end{itemize}

\subsubsection{Local Functions}
\begin{itemize}
\item \texttt{calculateResources()}: Computes resource changes resulting from a move.
\item \texttt{updateScores()}: Updates player score values based on new state.
\item \texttt{checkVictory()}: Determines whether any player meets winning conditions.
\end{itemize}

%%%%%%%%%%%%%%%%%%%%%%%%%%%%%%%%%%%%%%%%%%%%%%
\section{MIS of Training Pipeline} \label{M5}

\subsection{Module}
Training Pipeline (M5)

\subsection{Uses}

\begin{itemize}
    \item \textbf{M4}: Uses the access programs 
    \texttt{updateState}, \texttt{getState}, and \texttt{validateMove}.
    \item \textbf{AIMove}: Encoded AI-selected action.
    \item \textbf{GameStateData}: Abstract representation of a full game state.
    \item \textbf{Reward}: A real-valued score, $\mathbb{R}$.
    \item \textbf{StepError, RenderError}: Exceptions raised on invalid steps or rendering failures.
\end{itemize}

\subsection{Syntax}

\subsubsection{Exported Constants}
\begin{itemize}
\item \texttt{MAX\_TURNS}: $\mathbb{Z}^{+}$ - maximum number of turns in an episode
\item \texttt{REWARD\_SCALE}: $\mathbb{R}$ - scaling factor applied to reward calculation
\end{itemize}

\subsubsection{Exported Access Programs}
\begin{center}
\begin{tabular}{p{5cm} p{4cm} p{4cm} p{3cm}}
\hline
\textbf{Name} & \textbf{In} & \textbf{Out} & \textbf{Exceptions} \\
\hline
step & action:AIMove & (GameStateData, $\mathbb{R}$) & StepError \\
reset & - & GameStateData & - \\
render & - & - & RenderError \\
\hline
\end{tabular}
\end{center}

\subsection{Semantics}

\subsubsection{State Variables}
\begin{itemize}
\item \texttt{simulatedState}: internal instance of M4 \\
      Maintains the environment's evolving game state.

\item \texttt{turnCount}: $\mathbb{Z}_{\ge 0}$ \\
      Number of turns elapsed in the current episode.
\end{itemize}

\subsubsection{Environment Variables}
\begin{itemize}
\item \texttt{display}: Rendering surface or graphical context used to visualize simulatedState.
\end{itemize}

\subsubsection{Assumptions}
\begin{itemize}
    \item All inputs of type \texttt{AIMove} are syntactically valid.
    \item The environment is responsible for handling opponent moves internally.
\end{itemize}

\subsubsection{Access Routine Semantics}

\noindent step(action):
\begin{itemize}
\item transition:
The state change is governed by the rule:
\[
(
  \texttt{validateMove(action)} = \top \Rightarrow \text{apply actions}
|
  \texttt{validateMove(action)} = \bot \Rightarrow \text{-}
)
\]
\quad The "apply actions" transition includes the following sequential operations:
\begin{enumerate}
    \item Apply $\texttt{action}$ to $\texttt{simulatedState}$ via $\texttt{updateState(action)}$.
    \item Apply $\texttt{applyOpponentLogic()}$ to simulate adversarial moves.
    \item $\texttt{turnCount} := \texttt{turnCount} + 1$.
\end{enumerate}
\item output:  
      Returns the tuple $(\texttt{simulatedState.getState()}, \texttt{calculateReward()})$  
      of type $(\text{GameStateData}, \mathbb{R})$.
\item exception:  
      raises \texttt{StepError} if $\texttt{validateMove(action)} = \bot$.
\end{itemize}

\noindent reset():
\begin{itemize}
\item transition:  
      \texttt{simulatedState} := new initial game state (via M4).  
      turnCount := 0.
\item output:  
      Returns initial game state: GameStateData.
\item exception: -
\end{itemize}

\noindent render():
\begin{itemize}
\item transition:  
      Updates the \texttt{display} device to visualize the current \texttt{simulatedState}.
\item output: -
\item exception: raises \texttt{RenderError} if display rendering fails.
\end{itemize}

\subsubsection{Local Functions}
\begin{itemize}
\item \texttt{applyOpponentLogic()}: Generates and applies opponent actions based on heuristics or baseline policies.
\item \texttt{calculateReward()}: Returns a reward value in $\mathbb{R}$ based on changes in \texttt{simulatedState}, scaled by \texttt{REWARD\_SCALE}.
\end{itemize}

%%%%%%%%%%%%%%%%%%%%%%%%%%%%%%%%%%%%%%%%%%%%%%%
\section{MIS of AI Model Module} \label{M6}

\subsection{Module}
AI Model (M6)

\subsection{Uses}

\begin{itemize}
    \item \textbf{GameStateData}: Abstract representation of the game state.
    \item \textbf{AIMove}: Encoded move predicted by the model.
    \item \textbf{PredictionError}: Exception raised when the model cannot generate a valid move.
\end{itemize}

\subsection{Syntax}

\subsubsection{Exported Constants}
\begin{itemize}
\item \texttt{MODEL\_PATH}: $\Sigma^{*}$ - file path to stored model weights.
\end{itemize}

\subsubsection{Exported Access Programs}
\begin{center}
\begin{tabular}{p{5cm} p{4cm} p{4cm} p{3cm}}
\hline
\textbf{Name} & \textbf{In} & \textbf{Out} & \textbf{Exceptions} \\
\hline
decide & state:GameStateData & AIMove & PredictionError \\
getMoveConfidence & move:AIMove & $\mathbb{R}$ & - \\
explainMove & move:AIMove & $\Sigma^{*}$ & - \\
\hline
\end{tabular}
\end{center}

\subsection{Semantics}

\subsubsection{State Variables}
\begin{itemize}
\item \texttt{modelWeights}: internal representation of learned parameters  
      (vector or tensor values in $\mathbb{R}$ or $\mathbb{R}^{n}$).

\item \texttt{policyNetwork}: An abstract data structure representing the architecture and layers of the DRL model, initialized using \texttt{MODEL\_PATH}.

\item \texttt{lastPredictionCache}: dict mapping AIMove $\rightarrow \mathbb{R}$  
      storing confidence values for the most recent decision.
\end{itemize}

\subsubsection{Environment Variables}
None.

\subsubsection{Assumptions}
\begin{itemize}
    \item The input \texttt{GameStateData} conforms to the expected internal structure.
    \item \texttt{MODEL\_PATH} correctly references initialized modelWeights.
\end{itemize}

\subsubsection{Access Routine Semantics}

\noindent decide(state):
\begin{itemize}
\item transition:
    \begin{itemize}
        \item Converts $state$ using \texttt{preprocessState()}.
        \item Evaluates result  by calling \texttt{evaluatePolicy()}.
        \item Updates \texttt{lastPredictionCache} with move-confidence pairs.
    \end{itemize}

\item output:  
      Returns the optimal AIMove predicted by the model.

\item exception:  
      raises \texttt{PredictionError} if the model evaluation fails or produces an empty distribution.
\end{itemize}

\noindent getMoveConfidence(move):
\begin{itemize}
\item transition: -
\item output:  
      Returns $\mathbb{R}$ - the confidence score associated with $move$  
      retrieved from \texttt{lastPredictionCache}.
\item exception: -
\end{itemize}

\noindent explainMove(move):
\begin{itemize}
\item transition: -
\item output:  
      Returns $\Sigma^{*}$ - a human-readable textual explanation of the rationale  
      behind selecting $move$ (e.g., feature contributions or decision breakdown).
\item exception: -
\end{itemize}

\subsubsection{Local Functions}
\begin{itemize}
\item \texttt{preprocessState()}:
      Converts GameStateData into an abstract processed representation suitable  
      for model inference.

\item \texttt{postprocessOutput()}:
      Converts model policy output into a structured AIMove.
\item \texttt{evaluatePolicy()}: \texttt{ProcessedState} $\rightarrow$ \text{dist(AIMove)}
\begin{itemize}
      \item \text{Output}: Returns a probability distribution over possible moves by applying the \texttt{policyNetwork} data structure and \texttt{modelWeights} to the input state.
\end{itemize}
\end{itemize}

%%%%%%%%%%%%%%%%%%%%%%%%%%%%%%%%%%%%%%%%%%%%%%
\section{MIS of Game State Database Module} \label{M7}

\subsection{Module}
Game State Database (M7)

\subsection{Uses}

\begin{itemize}
    \item \textbf{GameStateData}: Abstract representation of a game state.
    \item \textbf{DBWriteError}: Exception raised on failed write operations.
    \item \textbf{DBReadError}: Exception raised on failed read or query operations.
    \item \textbf{seq(GameStateData)}: Ordered collection of game states.
\end{itemize}

\subsection{Syntax}

\subsubsection{Exported Constants}
\begin{itemize}
\item \texttt{DB\_PATH}: $\Sigma^{*}$ - file system path or database URI.
\item \texttt{MAX\_ENTRIES}: $\mathbb{Z}^{+}$ - maximum number of stored states.
\end{itemize}

\subsubsection{Exported Access Programs}
\begin{center}
\begin{tabular}{p{5cm} p{4cm} p{4cm} p{3cm}}
\hline
\textbf{Name} & \textbf{In} & \textbf{Out} & \textbf{Exceptions} \\
\hline
writeState & state:GameStateData & - & DBWriteError \\
readState & gameID:$\Sigma^{*}$ & GameStateData & DBReadError \\
queryHistory & playerID:$\Sigma^{*}$ & seq(GameStateData) & DBReadError \\
\hline
\end{tabular}
\end{center}

\subsection{Semantics}

\subsubsection{State Variables}
\begin{itemize}
\item \texttt{dbConnection}: Active connection handle to the database (abstract type).
\item \texttt{gameRecords}: dict mapping $\Sigma^{*}$ $\rightarrow$ GameStateData  
      (cache of recently accessed or frequently used game states).
\end{itemize}

\subsubsection{Environment Variables}
\begin{itemize}
\item \texttt{database}: External file system or DB server located at \texttt{DB\_PATH}.
\end{itemize}

\subsubsection{Assumptions}
\begin{itemize}
    \item The database at \texttt{DB\_PATH} is reachable and adheres to the expected schema.
    \item Identifiers such as \texttt{gameID} and \texttt{playerID} are valid $\Sigma^{*}$ strings.
\end{itemize}

\subsubsection{Access Routine Semantics}

\noindent writeState(state):
\begin{itemize}
\item transition:
    \begin{itemize}
        \item Converts $state$ to a serializable format via \texttt{serializeState()}.
        \item Writes serialized data to \texttt{database}.
        \item Updates \texttt{gameRecords} cache if applicable.
    \end{itemize}
\item output: -
\item exception: raises \texttt{DBWriteError} if the write operation fails.
\end{itemize}

\noindent readState(gameID):
\begin{itemize}
\item transition:
    \begin{itemize}
        \item If $gameID$ is cached, return entry from \texttt{gameRecords}.
        \item Else, load corresponding data from \texttt{database}.
        \item Deserialize via \texttt{deserializeState()}.
        \item Update \texttt{gameRecords}.
    \end{itemize}
\item output: GameStateData corresponding to the given $gameID$.
\item exception: raises \texttt{DBReadError} if no entry exists for $gameID$ or if read fails.
\end{itemize}

\noindent queryHistory(playerID):
\begin{itemize}
\item transition:
    \begin{itemize}
        \item Queries \texttt{database} for all entries linked to $playerID$.
        \item Converts results to seq(GameStateData) via \texttt{deserializeState()}.
    \end{itemize}
\item output: seq(GameStateData) containing all states associated with $playerID$.
\item exception: raises \texttt{DBReadError} if the query fails.
\end{itemize}

\subsubsection{Local Functions}
\begin{itemize}
\item \texttt{serializeState()}: Converts GameStateData into a serializable representation.
\item \texttt{deserializeState()}: Converts stored representation back to GameStateData.
\item \texttt{openConnection()}: Initializes \texttt{dbConnection}.
\item \texttt{closeConnection()}: Terminates \texttt{dbConnection}.
\end{itemize}

%%%%%%%%%%%%%%%%%%%%%%%%%%%%%%%%%%%%%%%%%%%%%%

\section{MIS of Image Queue Module (Phased Out)} \label{M8}

\subsection{Module}
Image Queue (M8).
In all other docs, M8 has been replaced by the Backend API Server module, and the image queue has been phased out.
However, in this MIS doc, the original design is kept as a legacy implementation to cover the formalization requirement
of the MIS rubric.
This formalization was created before we phased out the module, and it is retained here only for completeness.
The true M8 module, the Backend API Server, is defined immediately after this section.

\subsection{Uses}

\begin{itemize}
    \item \textbf{FrameData}: Abstract representation of an image frame from the camera or simulation.
    \item \textbf{QueueFullError}: Exception raised when attempting to enqueue a frame into a full queue.
    \item \textbf{QueueEmptyError}: Exception raised when attempting to dequeue or peek from an empty queue.
\end{itemize}

\subsection{Syntax}

\subsubsection{Exported Constants}
\begin{itemize}
\item \texttt{QUEUE\_SIZE} $\in \mathbb{Z}^{+}$ - maximum number of frames in the buffer.
\end{itemize}

\subsubsection{Exported Access Programs}
\begin{center}
\begin{tabular}{p{3cm} p{4cm} p{4cm} p{4cm}}
\hline
\textbf{Name} & \textbf{In} & \textbf{Out} & \textbf{Exceptions} \\
\hline
enqueue & frame:FrameData & - & QueueFullError \\
dequeue & - & FrameData & QueueEmptyError \\
peek & - & FrameData & QueueEmptyError \\
isFull & - & $\mathbb{B}$ & - \\
isEmpty & - & $\mathbb{B}$ & - \\
\hline
\end{tabular}
\end{center}

\subsection{Semantics}

\subsubsection{State Variables}
\begin{itemize}
\item \texttt{queueBuffer} : $[0..\text{QUEUE\_SIZE}-1] \rightarrow \text{FrameData} \cup \{\bot\}$ - circular buffer storing frames.
\item \texttt{head} : $\mathbb{Z}$ - index of the front of the queue.
\item \texttt{tail} : $\mathbb{Z}$ - index of the end of the queue.
\item \texttt{size} : $\mathbb{Z}$ - current number of frames in the queue.
\end{itemize}

\subsubsection{Environment Variables}
None.

\subsubsection{Assumptions}
\begin{itemize}
    \item Calling modules will check \texttt{isFull()} before \texttt{enqueue()}.
    \item Calling modules will check \texttt{isEmpty()} before \texttt{dequeue()} or \texttt{peek()}.
\end{itemize}

\subsubsection{Access Routine Semantics}

\noindent enqueue(frame: FrameData):
\begin{itemize}
\item transition:  
\[
\text{queueBuffer}[\text{tail}] := \text{frame}, \quad
\text{tail} := (\text{tail} + 1) \bmod \text{QUEUE\_SIZE}, \quad
\text{size} := \text{size} + 1
\]
\item output: -
\item exception: raises \texttt{QueueFullError} if $\text{isFull()} = \top$.
\end{itemize}

\noindent dequeue():
\begin{itemize}
\item transition: frame := queueBuffer[head]; queueBuffer[head] := $\bot$; head := (head+1) mod QUEUE\_SIZE; size := size-1
\item output: frame
\item exception: QueueEmptyError if isEmpty() = $\top$
\end{itemize}

\noindent peek():
\begin{itemize}
\item transition: -
\item output: $\text{queueBuffer}[\text{head}]$
\item exception: raises \texttt{QueueEmptyError} if $\text{isEmpty()} = \top$.
\end{itemize}

\noindent isFull():
\begin{itemize}
\item transition: -
\item output: $\top$ if $\text{size} = \text{QUEUE\_SIZE}$, $\bot$ otherwise
\item exception: -
\end{itemize}

\noindent isEmpty():
\begin{itemize}
\item transition: -
\item output: $\top$ if $\text{size} = 0$, $\bot$ otherwise
\item exception: -
\end{itemize}

\noindent resetQueue():
\begin{itemize}
\item transition:
\[
\text{head} := 0, \quad \text{tail} := 0, \quad \text{size} := 0
\]
\item output: -
\item exception: -
\end{itemize}

\subsubsection{Local Functions}
None.

\subsubsection{Formalization}

Let the state of this module be the tuple
\[
S = (queueBuffer, head, tail, size).
\]

\noindent \textbf{State Invariants} (must hold in all reachable states):

\begin{itemize}
    \item $0 \le size \le QUEUE\_SIZE$
    \item $head \in [0..QUEUE\_SIZE-1] \wedge tail \in [0..QUEUE\_SIZE-1]$
    \item $isEmpty() = \top \iff size = 0$
    \item $isFull() = \top \iff size = QUEUE\_SIZE$
\end{itemize}



\paragraph{enqueue(frame: FrameData)}
\textit{Defined when $isFull() = \bot$}
\begin{align*}
    &head' = head \\
    &queueBuffer'[tail] = frame \wedge tail' = (tail+1) \bmod QUEUE\_SIZE \wedge size' = size + 1 \\
    &\forall i \in [0..QUEUE\_SIZE-1], i \ne tail \Rightarrow queueBuffer'[i] = queueBuffer[i]
\end{align*}

\paragraph{dequeue()}
\textit{Defined when $isEmpty() = \bot$, output is $frame$}
\begin{align*}
    &frame = queueBuffer[head] \\
    &tail' = tail \\
    &queueBuffer'[head] = \bot \wedge head' = (head+1) \bmod QUEUE\_SIZE \wedge size' = size - 1 \\
    &\forall i \in [0..QUEUE\_SIZE-1], i \ne head \Rightarrow queueBuffer'[i] = queueBuffer[i]
\end{align*}

\paragraph{peek()}
\textit{Defined when $isEmpty() = \bot$, output is $frame$}
\begin{align*}
    &frame = queueBuffer[head] \\
    &queueBuffer' = queueBuffer \wedge head' = head \wedge tail' = tail \wedge size' = size
\end{align*}

\paragraph{resetQueue()}
\begin{align*}
    &head' = 0 \wedge tail' = 0 \wedge size' = 0 \wedge queueBuffer' = queueBuffer
\end{align*}


\section{MIS of Backend API Server Module} \label{M8}

\subsection{Module}
Backend API Server (M8)

\subsection{Uses}
\begin{itemize}
    \item \textbf{GameStateData}: Serialized game-state payloads.
    \item \textbf{MoveData}: Request or response action payload.
    \item \textbf{seq(BotMeta)}: Bot list returned to the frontend.
    \item \textbf{ExplanationText}: Natural-language explanation payload.
    \item \textbf{ProbabilityMap}: MCTS analysis payload.
    \item \textbf{M4.applyAction(), M4.getState(), M4.getActionContext()}.
    \item \textbf{M6.decide()}.
    \item \textbf{M7.readState(), M7.writeState(), M7.readLeagueMembers()}.
    \item \textbf{M9.listMembers()}.
    \item \textbf{M11.explainMove()}.
    \item \textbf{APIRequestError}: Raised when a request is invalid or cannot be fulfilled.
\end{itemize}

\subsection{Syntax}

\subsubsection{Exported Constants}
\begin{itemize}
\item \texttt{API\_PREFIX}: $\Sigma^{*}$ = \texttt{/api}
\end{itemize}

\subsubsection{Exported Access Programs}
\begin{center}
\begin{tabular}{p{5cm} p{4cm} p{4cm} p{3cm}}
\hline
\textbf{Name} & \textbf{In} & \textbf{Out} & \textbf{Exceptions} \\
\hline
listBots & - & seq(BotMeta) & APIRequestError \\
createGame & players:seq(PlayerKey) & $\Sigma^*$ & APIRequestError \\
getState & gameID:$\Sigma^*$, stateIndex:$\mathbb{Z} \cup \{\bot\}$ & GameStateData & APIRequestError \\
postAction & gameID:$\Sigma^*$, action:MoveData $\cup \{\bot\}$ & GameStateData & APIRequestError \\
getMCTSAnalysis & gameID:$\Sigma^*$, stateIndex:$\mathbb{Z} \cup \{\bot\}$ & ProbabilityMap & APIRequestError \\
getMoveExplanation & gameID:$\Sigma^*$, moveIndex:$\mathbb{N}$ & ExplanationText & APIRequestError \\
\hline
\end{tabular}
\end{center}

\subsection{Semantics}

\subsubsection{State Variables}
\begin{itemize}
\item \texttt{activeExplanationGameID}: $\Sigma^* \cup \{\bot\}$
\item \texttt{routeTable}: abstract mapping from HTTP route to handler
\end{itemize}

\subsubsection{Environment Variables}
\begin{itemize}
\item \texttt{httpRequest}: incoming request supplied by the web server runtime
\item \texttt{httpResponse}: outgoing response written to the client
\end{itemize}

\subsubsection{Assumptions}
\begin{itemize}
    \item Requests are transmitted as HTTP requests with JSON payloads where required.
    \item The backend is deployed with a valid database configuration and CORS-enabled web server environment.
\end{itemize}

\subsubsection{Access Routine Semantics}

\noindent listBots():
\begin{itemize}
\item transition: queries stored or configured league members
\item output: returns a list of bot metadata for UI selection
\item exception: raises \texttt{APIRequestError} if retrieval fails
\end{itemize}

\noindent createGame(players):
\begin{itemize}
\item transition: initializes a new game instance and persists its initial state
\item output: returns a generated game identifier
\item exception: raises \texttt{APIRequestError} if the request is malformed or initialization fails
\end{itemize}

\noindent getState(gameID, stateIndex):
\begin{itemize}
\item transition: -
\item output: returns latest or indexed historical game state
\item exception: raises \texttt{APIRequestError} if the requested state is unavailable
\end{itemize}

\noindent postAction(gameID, action):
\begin{itemize}
\item transition: applies a human move if \texttt{action} is supplied; otherwise advances the game with a bot move
\item output: returns the resulting game state
\item exception: raises \texttt{APIRequestError} if the request is invalid or the move cannot be applied
\end{itemize}

\noindent getMCTSAnalysis(gameID, stateIndex):
\begin{itemize}
\item transition: computes or retrieves win-probability analysis for the selected state
\item output: returns a probability map
\item exception: raises \texttt{APIRequestError} if analysis cannot be produced
\end{itemize}

\noindent getMoveExplanation(gameID, moveIndex):
\begin{itemize}
\item transition: gathers explanation context and delegates explanation generation to M11
\item output: returns explanation text
\item exception: raises \texttt{APIRequestError} if no explanation context exists or generation fails
\end{itemize}

\subsubsection{Local Functions}
\begin{itemize}
\item \texttt{parseStateIndex()}: validates replay-state selectors
\item \texttt{serializeResponse()}: converts backend results to JSON
\item \texttt{validateRequestBody()}: validates request shape before dispatch
\end{itemize}



\section{MIS of Elo League System Module} \label{M9}

\subsection{Module}
Elo League System (M9)

\subsection{Uses}
\begin{itemize}
    \item \textbf{LeagueEntry}: Stored rating record for a bot or model snapshot.
    \item \textbf{MatchResult}: Structured representation of a completed match outcome.
    \item \textbf{PlayerKey}: Identifier for a selectable bot.
    \item \textbf{seq(BotMeta)}: UI-visible league members.
    \item \textbf{M7.writeModelMetadata(), M7.readLeagueMembers()}.
    \item \textbf{LeagueError}: Raised when league operations fail.
\end{itemize}

\subsection{Syntax}

\subsubsection{Exported Constants}
\begin{itemize}
\item \texttt{MAX\_LEAGUE\_SIZE}: $\mathbb{N}$ = 125
\item \texttt{DEFAULT\_ELO}: $\mathbb{Z}$ = 1000
\end{itemize}

\subsubsection{Exported Access Programs}
\begin{center}
\begin{tabular}{p{5cm} p{4cm} p{4cm} p{3cm}}
\hline
\textbf{Name} & \textbf{In} & \textbf{Out} & \textbf{Exceptions} \\
\hline
registerMember & modelID:$\Sigma^*$ & PlayerKey & LeagueError \\
updateRatings & result:MatchResult & - & LeagueError \\
listMembers & - & seq(BotMeta) & LeagueError \\
selectOpponent & - & PlayerKey & LeagueError \\
pruneLeague & - & - & LeagueError \\
\hline
\end{tabular}
\end{center}

\subsection{Semantics}

\subsubsection{State Variables}
\begin{itemize}
\item \texttt{leagueTable}: seq(LeagueEntry)
\item \texttt{samplingWeights}: dict mapping PlayerKey $\rightarrow \mathbb{R}$
\end{itemize}

\subsubsection{Environment Variables}
None.

\subsubsection{Assumptions}
\begin{itemize}
    \item Match outcomes provided to the league are valid and complete.
    \item Stored model snapshots referenced by league entries are loadable.
\end{itemize}

\subsubsection{Access Routine Semantics}

\noindent registerMember(modelID):
\begin{itemize}
\item transition: adds a new league entry for the model snapshot
\item output: returns a selectable player key
\item exception: raises \texttt{LeagueError} if registration fails
\end{itemize}

\noindent updateRatings(result):
\begin{itemize}
\item transition: updates member Elo values according to the supplied match result
\item output: -
\item exception: raises \texttt{LeagueError} if ratings cannot be updated
\end{itemize}

\noindent listMembers():
\begin{itemize}
\item transition: -
\item output: returns league members in UI-consumable form
\item exception: raises \texttt{LeagueError} if retrieval fails
\end{itemize}

\noindent selectOpponent():
\begin{itemize}
\item transition: samples or chooses an opponent according to league policy
\item output: returns a player key
\item exception: raises \texttt{LeagueError} if no valid opponent is available
\end{itemize}

\noindent pruneLeague():
\begin{itemize}
\item transition: removes or archives low-priority members to keep the league bounded
\item output: -
\item exception: raises \texttt{LeagueError} if pruning fails
\end{itemize}

\subsubsection{Local Functions}
\begin{itemize}
\item \texttt{computeEloDelta()}: calculates rating change from a result
\item \texttt{normalizeWeights()}: updates sampling probabilities
\item \texttt{rankMembers()}: orders league entries for pruning or display
\end{itemize}



\section{MIS of Curriculum Learning Manager Module} \label{M10}

\subsection{Module}
Curriculum Learning Manager (M10)

\subsection{Uses}
\begin{itemize}
    \item \textbf{PhaseConfig}: Structured description of a curriculum phase.
    \item \textbf{TrainingMetrics}: Metrics used to decide advancement.
    \item \textbf{M9.selectOpponent()}.
    \item \textbf{CurriculumError}: Raised when curriculum state is invalid.
\end{itemize}

\subsection{Syntax}

\subsubsection{Exported Constants}
\begin{itemize}
\item \texttt{MIN\_PHASE\_COUNT}: $\mathbb{N}$
\end{itemize}

\subsubsection{Exported Access Programs}
\begin{center}
\begin{tabular}{p{5cm} p{4cm} p{4cm} p{3cm}}
\hline
\textbf{Name} & \textbf{In} & \textbf{Out} & \textbf{Exceptions} \\
\hline
getCurrentPhase & runID:$\Sigma^*$ & PhaseConfig & CurriculumError \\
evaluateAdvance & metrics:TrainingMetrics & $\mathbb{B}$ & CurriculumError \\
advancePhase & runID:$\Sigma^*$ & PhaseConfig & CurriculumError \\
logPhaseTransition & runID:$\Sigma^*$, phase:PhaseConfig & - & CurriculumError \\
\hline
\end{tabular}
\end{center}

\subsection{Semantics}

\subsubsection{State Variables}
\begin{itemize}
\item \texttt{phaseTable}: seq(PhaseConfig)
\item \texttt{activePhase}: dict mapping run identifiers to phase index
\end{itemize}

\subsubsection{Environment Variables}
None.

\subsubsection{Assumptions}
\begin{itemize}
    \item Curriculum phases are defined in ascending order of complexity.
    \item Training metrics supplied to this module are internally consistent.
\end{itemize}

\subsubsection{Access Routine Semantics}

\noindent getCurrentPhase(runID):
\begin{itemize}
\item transition: -
\item output: returns the active phase configuration for the run
\item exception: raises \texttt{CurriculumError} if the run has no active phase
\end{itemize}

\noindent evaluateAdvance(metrics):
\begin{itemize}
\item transition: -
\item output: returns $\top$ iff advancement thresholds are met
\item exception: raises \texttt{CurriculumError} if metrics are invalid
\end{itemize}

\noindent advancePhase(runID):
\begin{itemize}
\item transition: increments or cycles the active phase for the run according to policy
\item output: returns the new phase configuration
\item exception: raises \texttt{CurriculumError} if no valid next phase exists
\end{itemize}

\noindent logPhaseTransition(runID, phase):
\begin{itemize}
\item transition: records that the run entered \texttt{phase}
\item output: -
\item exception: raises \texttt{CurriculumError} if logging fails
\end{itemize}

\subsubsection{Local Functions}
\begin{itemize}
\item \texttt{meetsThresholds()}: checks advancement criteria
\item \texttt{buildPhaseRules()}: materializes rule variations for a phase
\item \texttt{cyclePolicy()}: applies fallback/cycling logic
\end{itemize}


\section{MIS of Explanation Pipeline Module} \label{M11}

\subsection{Module}
Explanation Pipeline (M11)

\subsection{Uses}
\begin{itemize}
    \item \textbf{GameStateData}: State used as explanation context.
    \item \textbf{MoveData}: Selected move to be explained.
    \item \textbf{ExplanationText}: Natural-language explanation output.
    \item \textbf{dict}: Structured prompt/context packet.
    \item \textbf{M4.getActionContext()}.
    \item \textbf{ExplanationError}: Raised when explanation generation fails.
\end{itemize}

\subsection{Syntax}

\subsubsection{Exported Constants}
\begin{itemize}
\item \texttt{MAX\_CONTEXT\_ACTIONS}: $\mathbb{N}$
\end{itemize}

\subsubsection{Exported Access Programs}
\begin{center}
\begin{tabular}{p{5cm} p{4cm} p{4cm} p{3cm}}
\hline
\textbf{Name} & \textbf{In} & \textbf{Out} & \textbf{Exceptions} \\
\hline
buildContextPacket & gameID:$\Sigma^*$, moveIndex:$\mathbb{N}$ & dict & ExplanationError \\
explainMove & gameID:$\Sigma^*$, moveIndex:$\mathbb{N}$ & ExplanationText & ExplanationError \\
validateExplanation & text:ExplanationText & $\mathbb{B}$ & - \\
\hline
\end{tabular}
\end{center}

\subsection{Semantics}

\subsubsection{State Variables}
\begin{itemize}
\item \texttt{lastContextPacket}: dict $\cup \{\bot\}$
\item \texttt{lastExplanation}: ExplanationText $\cup \{\bot\}$
\end{itemize}

\subsubsection{Environment Variables}
\begin{itemize}
\item \texttt{llmService}: external large-language-model service
\end{itemize}

\subsubsection{Assumptions}
\begin{itemize}
    \item Sufficient replay context exists for the selected move.
    \item The external LLM service may fail or rate-limit requests.
\end{itemize}

\subsubsection{Access Routine Semantics}

\noindent buildContextPacket(gameID, moveIndex):
\begin{itemize}
\item transition: gathers structured move and state context from the game-state manager and stores it in \texttt{lastContextPacket}
\item output: returns the generated context dictionary
\item exception: raises \texttt{ExplanationError} if context cannot be built
\end{itemize}

\noindent explainMove(gameID, moveIndex):
\begin{itemize}
\item transition: builds context if needed, submits the explanation request to the external LLM service, and stores the result in \texttt{lastExplanation}
\item output: returns natural-language explanation text
\item exception: raises \texttt{ExplanationError} if generation fails or no explanation is available
\end{itemize}

\noindent validateExplanation(text):
\begin{itemize}
\item transition: -
\item output: returns $\top$ iff the explanation is non-empty and meets formatting constraints
\item exception: -
\end{itemize}

\subsubsection{Local Functions}
\begin{itemize}
\item \texttt{formatPrompt()}: converts context into an LLM-ready prompt
\item \texttt{sanitizeResponse()}: removes unusable or malformed output
\item \texttt{fallbackErrorText()}: generates safe user-facing failure text
\end{itemize}

%%%%%%%%%END OF MODULES ADDED%%%%%%%%%%%%%%%%%%%%
\newpage

\bibliographystyle {plainnat}
\bibliography {../../../refs/References}

\newpage{}


\section*{Appendix - Reflection}


%\input{../../Reflection.text}

\begin{enumerate}
  \item What went well while writing this deliverable? 
   
  [Group] Our team collaborated effectively, dividing sections based on
          what each member was most familiar with. Rebecca and Matthew
          focused on the MIS and MG since they had the most experience with
          these documents from before. While Sunny and Jake focused on the Software
          during this time.
  \item What pain points did you experience during this deliverable, and how  
    did you resolve them?

  [Group] One challenge we faced was ensuring consistency across the documents
          and ensuring our modules were decoupled enough. We resolved this by
          having multiple review sessions where we cross-checked each other's
          sections and provided feedback. This iterative process helped us
          create modules that were well defined and cohesive.
  
  
  \item Which of your design decisions stemmed from speaking to your client(s)
  or a proxy (e.g. your peers, stakeholders, potential users)? For those that
  were not, why, and where did they come from?
 
  [Group] A lot of our design decisions were influenced by discussions with
          our supervisor who provided insights into the software architecture
          of RL environments. We also consulted with peers who had experience
          in computer vision to refine our CV module design. Some decisions,
          like the choice of data structures, were based on our own research
          and understanding of best practices in software design.

  \item While creating the design doc, what parts of your other documents (e.g.
  requirements, hazard analysis, etc), it any, needed to be changed, and why?
  unlimited resources, what could you do to make the project better? (LO\_ProbSolutions)

   [Group]   While developing the design document, several updates were 
   required in our earlier project artifacts, particularly the hazard 
   analysis and system requirements.

The most significant change came from refining our system architecture. 
During hazard analysis we originally treated the digital twin as a 
simple board state representation. However, while writing the design 
document we clarified that the digital twin must represent not only 
the board layout, but also the full game state, player actions, and 
system input and output interactions. Because of this expanded definition, 
multiple hazards had to be revised. Risks that were previously limited 
to incorrect board modeling were broadened to include incorrect state 
 synchronization, input misinterpretation, and output decision errors. 
 This led to updated mitigation strategies, such as stronger state
 validation checks and clearer interface boundaries between the 
 AI agent and the game environment.

Our requirements document also changed. Some functional requirements 
were rewritten to reflect that the system must track and process full 
game context. Interface requirements became more detailed to define 
how the model, simulation environment, and data storage components 
communicate.

These revisions improved consistency across documents and ensured 
the design accurately reflected the true scope and risks of the system.

With additional resources, we would rent more GPUs to run parallel 
training jobs, allowing faster experimentation and more extensive model 
tuning. We would also pay for access to online Catan game history 
databases to expand our training data and improve model performance 
and generalization.


  \item Give a brief overview of other design solutions you considered.  What
  are the benefits and tradeoffs of those other designs compared with the chosen
  design?  From all the potential options, why did you select the documented design?
  (LO\_Explores)

[Group] We explored several alternative design approaches before selecting 
our final solution. One option was a simpler rule based system that relied 
on handcrafted heuristics rather than a learning model. The main benefit 
of this approach was predictability and lower computational cost, as it 
would not require long training times or large datasets. However, the 
tradeoff was limited adaptability and weaker performance in complex or 
unexpected game situations, since the system could only act within 
predefined rules.


  \item (After you have implemented another team's module, which means this
  isn't filled in until after the original deadline). What did you learn by
  implementing another team's module? Were all the details you needed in the
  documentation, or did you need to make assumptions, or ask the other team
  questions?  If your team also had another team implement one of your modules,
  what was this experience like?  Are there things in your documentation you
  could have changed to make the process go more smoothly for when an
  ``outsider'' completes some of the implementation?


  [Group] Implementing the assigned team's module was a significant learning experience in API design.
          While their src/reporting\_poc\_README.md outlined the general API surface, it lacked concrete request/response payloads, formal error codes, and a runnable test.
          This documentation gap forced us to move forward based on assumptions and manual follow-up questions rather than clear technical specifications.
          We forked their repository and opened a pull request to merge our code, intentionally seeking their review comments to validate our integration.
          We noted that because a PR-driven workflow wasn't used by the team, we missed out on automated CI/CD feedback, which made validating our integration assumptions more difficult.
          Moving forward, we’ve realized that truly exceptional modules must include explicit authorization contracts, a short integration checklist, and a tiny runnable test or CI workflow.
          This ensures that any implementer has a clear source of truth to reproduce expected behavior without needing to ask targeted questions.


\end{enumerate}

\end{document}